\documentclass[11pt]{article}

\usepackage[margin=1in]{geometry}
\usepackage{amsmath,amssymb,mathtools}
\usepackage{bm}
\usepackage{enumitem}
\usepackage{hyperref}

\newcommand{\ii}{\mathrm{i}}
\newcommand{\e}{\mathrm{e}}
\newcommand{\dd}{\mathrm{d}}
\newcommand{\1}{\mathbb{1}}
\newcommand{\bI}{\mathbf{I}}
\newcommand{\bh}{\mathbf{h}}
\newcommand{\diag}{\operatorname{diag}}
\newcommand{\cond}{\operatorname{cond}}
\newcommand{\abs}[1]{\left|#1\right|}
\newcommand{\br}[1]{\left(#1\right)}
\newcommand{\sq}[1]{\left[#1\right]}
\newcommand{\set}[1]{\left\{#1\right\}}

\title{Stable Numerical Evaluation of the Free Graded-Replica Green Kernel}
\author{}
\date{}

\begin{document}
\maketitle

\section{Purpose}

This note only discusses the numerical evaluation of the free
graded-replica Green function.  The path-integral construction, the graded
swap, and the perturbation theory are not repeated here.  The goal is to
start from the exact matrix formula for the free sector Green function and
rewrite it so that the large factors in
\(\exp(\pm \beta \bh)\) are never formed as large floating-point numbers.

The issue is the following.  The one-particle hopping Hamiltonian \(\bh\)
can have negative eigenvalues.  Then \(\e^{-\beta \bh}\) may contain large
eigenvalues.  In the exact Green function these large factors are cancelled
by the inverse matrix
\[
    \sq{\bI+\Omega_m^{-1}\e^{-\beta\bh}}^{-1},
\]
but a direct implementation forms the large matrix first and relies on the
linear algebra to cancel it later.  This is numerically fragile.  The stable
form below performs the cancellation algebraically before exponentiation.

\section{Original Sector Formula}

Let \(A\) be the entanglement subregion and \(B\) its complement.  On the
one-particle Hilbert space define the projectors
\[
    P_A+P_B=\bI,
    \qquad
    P_A P_B=0.
\]
For \(n\) replicas the graded momenta are
\[
    q_m=\frac{(2m+1)\pi}{n},
    \qquad
    m=0,\ldots,n-1.
\]
The sector gluing matrix is
\[
    \Omega_m = P_B-\e^{\ii q_m}P_A.
\]
Equivalently,
\[
    \Omega_m^{-1}=P_B-\e^{-\ii q_m}P_A.
\]
Both \(\Omega_m\) and \(\Omega_m^{-1}\) are unitary, because they are
diagonal phase matrices in the real-space \(A/B\) decomposition.

The free sector Green function satisfies
\[
    \br{\partial_\tau+\bh}\bar G_m(\tau,\tau')
    =
    \delta(\tau-\tau')\bI
\]
with sector boundary condition
\[
    \bar G_m(\beta^-,\tau')
    =
    -\Omega_m \bar G_m(0^+,\tau').
\]
Solving the equation of motion gives the exact piecewise formula
\[
    \bar G_m(\tau,\tau')
    =
    \begin{cases}
    \e^{-\tau\bh} A_m \e^{\tau'\bh},
    & \tau>\tau',
    \\[4pt]
    \e^{-\tau\bh} \br{A_m-\bI}\e^{\tau'\bh},
    & \tau<\tau',
    \end{cases}
\]
where
\[
    A_m =
    \sq{\bI+\Omega_m^{-1}\e^{-\beta\bh}}^{-1}.
\]

This formula is mathematically correct, but it is not the best numerical
formula when \(\beta \abs{\epsilon_a}\) is large.

\section{Diagonalizing the Hopping Hamiltonian}

Assume that \(\bh\) is real symmetric.  Write
\[
    \bh = V\Lambda V^T,
    \qquad
    V^T V=\bI,
    \qquad
    \Lambda_{ab}=\epsilon_a\delta_{ab}.
\]
Define the gluing matrix in the eigenbasis of \(\bh\):
\[
    S_m = V^T\Omega_m^{-1}V.
\]
Since \(V\) is real orthogonal and \(\Omega_m^{-1}\) is unitary,
\[
    S_m^\dagger S_m=\bI.
\]
Thus \(S_m\) is unitary.  In particular,
\[
    \abs{(S_m)_{ab}}\le 1.
\]

Let
\[
    D = \e^{\beta\Lambda}
    =
    \diag\br{\e^{\beta\epsilon_1},\ldots,\e^{\beta\epsilon_N}},
    \qquad
    D^{-1}=\e^{-\beta\Lambda}.
\]
In the eigenbasis,
\[
    \mathcal A_m
    \equiv
    V^T A_m V
    =
    \sq{\bI+S_mD^{-1}}^{-1}.
\]
Now factor \(D^{-1}\) from the right:
\[
    \bI+S_mD^{-1}
    =
    (D+S_m)D^{-1}.
\]
Therefore
\[
    \mathcal A_m
    =
    D(D+S_m)^{-1}.
\]
No commutativity between \(D\) and \(S_m\) is assumed.  The order in the
last formula is essential.

\section{Eigenbasis Kernels}

Introduce dimensionless times
\[
    u=\frac{\tau}{\beta},
    \qquad
    v=\frac{\tau'}{\beta},
    \qquad
    0\le u,v\le 1.
\]
The site-basis Green function is reconstructed from an eigenbasis kernel:
\[
    \bar G_m(\tau,\tau') = V K_m(\tau,\tau') V^T.
\]

\subsection{Upper branch}

For \(\tau>\tau'\), equivalently \(u>v\),
\[
    K_m^>(u,v)
    =
    D^{-u}\mathcal A_m D^v.
\]
Using \(\mathcal A_m=D(D+S_m)^{-1}\), this becomes
\[
    \boxed{
    K_m^>(u,v)
    =
    D^{1-u}(D+S_m)^{-1}D^v.
    }
\]
The raw matrix element form is
\[
    \br{K_m^>}_{ab}
    =
    \e^{(1-u)\beta\epsilon_a}
    \sq{(D+S_m)^{-1}}_{ab}
    \e^{v\beta\epsilon_b}.
\]
This already shows the desired cancellation structure.  The exponential
from \(D\), which came from inverting the original \(D^{-1}\), has been
combined with the external factors \(\e^{-\tau\Lambda}\) and
\(\e^{\tau'\Lambda}\).

\subsection{Lower branch}

For \(\tau<\tau'\), equivalently \(u<v\),
\[
    K_m^<(u,v)
    =
    D^{-u}\br{\mathcal A_m-\bI}D^v.
\]
Use
\[
    \mathcal A_m-\bI
    =
    -\sq{\bI+D S_m^{-1}}^{-1}.
\]
One way to see this without assuming any commutativity is
\[
    \mathcal A_m
    =
    D(D+S_m)^{-1},
\]
so
\[
    \mathcal A_m-\bI
    =
    D(D+S_m)^{-1}-(D+S_m)(D+S_m)^{-1}
    =
    -S_m(D+S_m)^{-1}.
\]
On the other hand,
\[
    \bI+D S_m^{-1}
    =
    (D+S_m)S_m^{-1},
\]
and therefore
\[
    \sq{\bI+D S_m^{-1}}^{-1}
    =
    S_m(D+S_m)^{-1}.
\]
Since
\[
    \bI+D S_m^{-1}
    =
    D\br{D^{-1}+S_m^{-1}},
\]
define
\[
    H = D^{-1}=\e^{-\beta\Lambda}.
\]
Then
\[
    \boxed{
    K_m^<(u,v)
    =
    -H^u\br{H+S_m^{-1}}^{-1}H^{1-v}.
    }
\]
Again, no commutativity is assumed.  Also,
\[
    S_m^{-1}=S_m^\dagger,
\]
so the lower branch does not require a separate matrix inverse for \(S_m\).

\section{Scalar Check: \(S_m=\bI\)}

If there is no mode mixing, \(S_m=\bI\), the upper branch reduces to
\[
    K_a^>(u,v)
    =
    \frac{D_a^{1-u+v}}{1+D_a}
    =
    \frac{\e^{(1-u+v)x_a}}{1+\e^{x_a}},
    \qquad
    x_a=\beta\epsilon_a,
\]
with \(0<1-u+v<1\).  This is a scalar Fermi kernel.  It is bounded:
\[
    x_a\gg 1:
    \quad
    K_a^>\sim \e^{-(u-v)x_a},
    \qquad
    x_a\ll -1:
    \quad
    K_a^>\sim \e^{(1-u+v)x_a}.
\]
Both asymptotic forms are small or \(O(1)\), never exponentially large.

The lower branch becomes
\[
    K_a^<(u,v)
    =
    -\frac{H_a^{u+1-v}}{1+H_a}
    =
    -\frac{\e^{(u+1-v)y_a}}{1+\e^{y_a}},
    \qquad
    y_a=-x_a,
\]
with \(0<u+1-v<1\).  It is bounded for the same reason.

The matrix formulas below are the noncommuting version of this scalar
``manually cancel the large exponential'' evaluation.

\subsection{Relation to the usual scalar cutoff trick}

For a scalar Fermi factor one often evaluates
\[
    f_r(x)=\frac{\e^{rx}}{1+\e^x},
    \qquad
    0\le r\le 1,
\]
by splitting into large-positive and large-negative regimes:
\[
    x\gg 1:
    \quad
    f_r(x)
    =
    \frac{\e^{(r-1)x}}{1+\e^{-x}}
    \simeq
    \e^{(r-1)x},
\]
\[
    x\ll -1:
    \quad
    f_r(x)
    =
    \frac{\e^{rx}}{1+\e^x}
    \simeq
    \e^{rx}.
\]
This is stable because the large exponential has been cancelled before the
number is evaluated.  A practical implementation may introduce a cutoff
\(T\): for \(x>T\) use the first asymptotic form, for \(x<-T\) use the
second, and otherwise use the original expression.

The diagonal scaling used in this note is the exact no-cutoff version of the
same idea.  Define
\[
    c=\max(0,x).
\]
Then
\[
    1+\e^x
    =
    \e^c\br{\e^{-c}+\e^{x-c}},
\]
and hence
\[
    f_r(x)
    =
    \frac{\e^{rx-c}}{\e^{-c}+\e^{x-c}}.
\]
All exponents in this expression are non-positive.  If \(x\ge0\), then
\[
    f_r(x)
    =
    \frac{\e^{(r-1)x}}{1+\e^{-x}},
\]
while if \(x<0\), then
\[
    f_r(x)
    =
    \frac{\e^{rx}}{1+\e^x}.
\]
Thus the scaling reproduces the usual stable scalar branches exactly.  The
cutoff method is recovered only as a further approximation when
\(\e^{-x}\) or \(\e^x\) is below the desired tolerance.  The matrix scaling
below keeps the exact scaled denominator instead of dropping the small term.

\subsection{Scalar limit of the matrix scaling}

When \(S_m=\bI\), the upper scaled middle matrix is diagonal:
\[
    M_a^+
    =
    \e^{-\theta_+c_a^+}\e^{-(1-\theta_+)c_a^+}
    +
    \e^{x_a-c_a^+}
    =
    \e^{-c_a^+}+\e^{x_a-c_a^+}.
\]
It is independent of \(\theta_+\).  The left and right factors are
\[
    L_a^+
    =
    \e^{(1-u)x_a-(1-\theta_+)c_a^+},
    \qquad
    R_a^+
    =
    \e^{vx_a-\theta_+c_a^+}.
\]
Therefore
\[
    L_a^+(M_a^+)^{-1}R_a^+
    =
    \frac{
    \e^{(1-u+v)x_a-c_a^+}
    }{
    \e^{-c_a^+}+\e^{x_a-c_a^+}
    }
    =
    \frac{\e^{(1-u+v)x_a}}{1+\e^{x_a}}.
\]
This is exactly the scalar upper Fermi kernel.

The lower branch is the same statement with
\[
    y_a=-x_a,
    \qquad
    c_a^-=\max(0,y_a).
\]
When \(S_m=\bI\),
\[
    M_a^-=\e^{-c_a^-}+\e^{y_a-c_a^-},
\]
and
\[
    -L_a^-(M_a^-)^{-1}R_a^-
    =
    -
    \frac{\e^{(u+1-v)y_a}}{1+\e^{y_a}}.
\]
Multiplying numerator and denominator by \(\e^{x_a}\) gives the equivalent
form
\[
    -
    \frac{\e^{(v-u)x_a}}{1+\e^{x_a}},
\]
which is the scalar lower Fermi kernel.  Thus in the commuting scalar limit
the scaled matrix formula is not merely analogous to the cutoff trick; it is
the exact stable logistic/Fermi evaluation.

\section{Stable Diagonal Scaling: Upper Branch}

The upper branch is
\[
    K_m^> = D^{1-u}(D+S_m)^{-1}D^v,
    \qquad
    u>v.
\]
Let
\[
    x_a=\beta\epsilon_a,
    \qquad
    c_a^+=\max(0,x_a),
    \qquad
    C_+ = \diag\br{\e^{c_a^+}}.
\]
Thus
\[
    (C_+)_{aa}=\max(1,D_a).
\]
Choose any
\[
    \theta_+\in[v,u].
\]
Define the scaled middle matrix
\[
    \boxed{
    M_m^+
    =
    C_+^{-\theta_+}
    \br{D+S_m}
    C_+^{-(1-\theta_+)}.
    }
\]
Then
\[
    D+S_m
    =
    C_+^{\theta_+} M_m^+ C_+^{1-\theta_+},
\]
so
\[
    (D+S_m)^{-1}
    =
    C_+^{-(1-\theta_+)}
    (M_m^+)^{-1}
    C_+^{-\theta_+}.
\]
Therefore
\[
    \boxed{
    K_m^>
    =
    L_+ (M_m^+)^{-1} R_+,
    }
\]
with diagonal prefactors
\[
    L_+
    =
    D^{1-u}C_+^{-(1-\theta_+)},
    \qquad
    R_+
    =
    C_+^{-\theta_+}D^v.
\]

\subsection{Elementwise form used in code}

The stable implementation should never form \(D\) or \(C_+\) as large
intermediate numbers.  It should work with exponents:
\[
    \ell_a^+
    =
    (1-u)x_a-(1-\theta_+)c_a^+,
\]
\[
    r_b^+
    =
    vx_b-\theta_+ c_b^+.
\]
Then
\[
    (L_+)_{aa}=\e^{\ell_a^+},
    \qquad
    (R_+)_{bb}=\e^{r_b^+}.
\]
The choice \(\theta_+\in[v,u]\) makes both exponents non-positive.
Indeed, if \(x_a\ge 0\), then \(c_a^+=x_a\) and
\[
    \ell_a^+=(\theta_+-u)x_a\le 0.
\]
If \(x_a<0\), then \(c_a^+=0\) and
\[
    \ell_a^+=(1-u)x_a\le 0.
\]
Similarly, if \(x_b\ge 0\),
\[
    r_b^+=(v-\theta_+)x_b\le 0,
\]
while if \(x_b<0\),
\[
    r_b^+=vx_b\le 0.
\]
Thus \(0<(L_+)_{aa}\le 1\) and \(0<(R_+)_{bb}\le 1\).

The entries of \(M_m^+\) are
\[
    \boxed{
    (M_m^+)_{ab}
    =
    \e^{-\theta_+ c_a^+}
    (S_m)_{ab}
    \e^{-(1-\theta_+)c_b^+}
    +
    \delta_{ab}\e^{x_a-c_a^+}.
    }
\]
Every exponent here is non-positive:
\[
    -\theta_+c_a^+\le 0,
    \qquad
    -(1-\theta_+)c_b^+\le 0,
    \qquad
    x_a-c_a^+\le 0.
\]
Hence no overflow can occur in the construction of \(M_m^+\), \(L_+\), or
\(R_+\).

\subsection{Why the upper middle matrix is bounded}

Because \(S_m\) is unitary,
\[
    \abs{(S_m)_{ab}}\le 1.
\]
The mixing term satisfies
\[
    \abs{
    \e^{-\theta_+ c_a^+}
    (S_m)_{ab}
    \e^{-(1-\theta_+)c_b^+}
    }
    \le 1.
\]
The diagonal term satisfies
\[
    0<\e^{x_a-c_a^+}\le 1.
\]
Thus each entry obeys the simple bound
\[
    \abs{(M_m^+)_{ab}}\le 1+\delta_{ab}.
\]
The middle matrix itself contains no exponentially large entries.

The inverse \((M_m^+)^{-1}\) can still be large if \(M_m^+\) is nearly
singular.  That is a genuine conditioning or physical zero-mode issue, not
an artificial overflow caused by \(\e^{\beta\abs{\epsilon}}\).

\section{Stable Diagonal Scaling: Lower Branch}

The lower branch is
\[
    K_m^<
    =
    -H^u(H+S_m^{-1})^{-1}H^{1-v},
    \qquad
    H=D^{-1},
    \qquad
    u<v.
\]
Let
\[
    y_a=-x_a=-\beta\epsilon_a,
    \qquad
    c_a^-=\max(0,y_a),
    \qquad
    C_-=\diag\br{\e^{c_a^-}}.
\]
Thus
\[
    (C_-)_{aa}=\max(1,H_a).
\]
Choose any
\[
    \theta_-\in[1-v,1-u].
\]
Define
\[
    \boxed{
    M_m^-
    =
    C_-^{-\theta_-}
    \br{H+S_m^{-1}}
    C_-^{-(1-\theta_-)}.
    }
\]
Then
\[
    \boxed{
    K_m^<
    =
    -L_-(M_m^-)^{-1}R_-,
    }
\]
with
\[
    L_-
    =
    H^u C_-^{-(1-\theta_-)},
    \qquad
    R_-
    =
    C_-^{-\theta_-}H^{1-v}.
\]

\subsection{Elementwise form used in code}

Again, use only combined exponents:
\[
    \ell_a^-
    =
    uy_a-(1-\theta_-)c_a^-,
\]
\[
    r_b^-
    =
    (1-v)y_b-\theta_- c_b^-.
\]
Then
\[
    (L_-)_{aa}=\e^{\ell_a^-},
    \qquad
    (R_-)_{bb}=\e^{r_b^-}.
\]
The interval \(\theta_-\in[1-v,1-u]\) makes these exponents non-positive.
If \(y_a\ge 0\), then \(c_a^-=y_a\) and
\[
    \ell_a^-=(u+\theta_- -1)y_a\le 0,
\]
because \(\theta_-\le 1-u\).  If \(y_a<0\), then \(c_a^-=0\) and
\[
    \ell_a^-=uy_a\le 0.
\]
For the right factor, if \(y_b\ge 0\),
\[
    r_b^-=(1-v-\theta_-)y_b\le 0,
\]
because \(\theta_-\ge 1-v\).  If \(y_b<0\), then
\[
    r_b^-=(1-v)y_b\le 0.
\]
Thus \(0<(L_-)_{aa}\le 1\) and \(0<(R_-)_{bb}\le 1\).

The entries of \(M_m^-\) are
\[
    \boxed{
    (M_m^-)_{ab}
    =
    \e^{-\theta_- c_a^-}
    (S_m^{-1})_{ab}
    \e^{-(1-\theta_-)c_b^-}
    +
    \delta_{ab}\e^{y_a-c_a^-}.
    }
\]
Since \(S_m^{-1}=S_m^\dagger\) is unitary,
\[
    \abs{(S_m^{-1})_{ab}}\le 1.
\]
Also all exponents are non-positive:
\[
    -\theta_-c_a^-\le 0,
    \qquad
    -(1-\theta_-)c_b^-\le 0,
    \qquad
    y_a-c_a^-\le 0.
\]
Therefore
\[
    \abs{(M_m^-)_{ab}}\le 1+\delta_{ab}.
\]
The lower middle matrix also contains no exponentially large entries.

\section{Choice of the Splitting Parameter}

The parameters \(\theta_+\) and \(\theta_-\) are not physical.  They only
choose how the diagonal scaling is divided between the left and right side
of the linear solve.  A balanced default is
\[
    \theta_+ = \frac{u+v}{2},
    \qquad
    \tau>\tau',
\]
and
\[
    \theta_- = \frac{(1-v)+(1-u)}{2}
    =
    1-\frac{u+v}{2},
    \qquad
    \tau<\tau'.
\]
These choices lie in the required intervals:
\[
    \theta_+\in[v,u],
    \qquad
    \theta_-\in[1-v,1-u].
\]
They also distribute the small factors symmetrically around the middle
matrix, which usually improves conditioning of the scaled linear solve.

\section{Practical Numerical Algorithm}

The following is the intended numerical workflow in the eigenbasis.

\subsection*{Precomputation for each spin and sector}

\begin{enumerate}[leftmargin=2em]
    \item Diagonalize the real symmetric hopping matrix:
    \[
        \bh=V\Lambda V^T.
    \]
    Store the eigenvalues \(\epsilon_a\) and eigenvectors \(V_{ia}\).

    \item Build the site-basis diagonal phases of \(\Omega_m^{-1}\):
    \[
        \eta_i^{(m)}
        =
        \begin{cases}
        1, & i\in B,\\
        -\e^{-\ii q_m}, & i\in A.
        \end{cases}
    \]
    If the code uses the conjugate Fourier convention, the sign in the
    phase should follow the convention used in the original \(B_m\).
    The numerical method below only requires that the same
    \(\Omega_m^{-1}\) used in \(B_m=\bI+\Omega_m^{-1}\e^{-\beta\bh}\)
    is used here.

    \item Build
    \[
        (S_m)_{ab}
        =
        \sum_i V_{ia}\eta_i^{(m)}V_{ib}.
    \]
    This is a dense unitary matrix, but it contains no large or small
    exponentials.

    \item Store
    \[
        x_a=\beta\epsilon_a.
    \]
    Do not store \(D_a=\e^{x_a}\) as the primary object.
\end{enumerate}

\subsection*{Evaluation for \(\tau>\tau'\)}

\begin{enumerate}[leftmargin=2em]
    \item Compute
    \[
        u=\tau/\beta,
        \qquad
        v=\tau'/\beta,
        \qquad
        \theta_+=(u+v)/2.
    \]

    \item For each eigenmode \(a\), compute
    \[
        c_a^+=\max(0,x_a).
    \]

    \item Build the scaled middle matrix
    \[
        (M_m^+)_{ab}
        =
        \e^{-\theta_+ c_a^+}
        (S_m)_{ab}
        \e^{-(1-\theta_+)c_b^+}
        +
        \delta_{ab}\e^{x_a-c_a^+}.
    \]

    \item Build the diagonal left and right factors from combined exponents:
    \[
        (L_+)_{aa}
        =
        \exp\sq{(1-u)x_a-(1-\theta_+)c_a^+},
    \]
    \[
        (R_+)_{aa}
        =
        \exp\sq{vx_a-\theta_+c_a^+}.
    \]

    \item Solve the linear system
    \[
        M_m^+Y=R_+
    \]
    for \(Y\).  Do not explicitly form \((M_m^+)^{-1}\) unless this is only
    for diagnostic testing.

    \item Form
    \[
        K_m^> = L_+Y.
    \]

    \item Transform back to site basis:
    \[
        \bar G_m^>(\tau,\tau')=VK_m^>V^T.
    \]
\end{enumerate}

\subsection*{Evaluation for \(\tau<\tau'\)}

\begin{enumerate}[leftmargin=2em]
    \item Compute
    \[
        u=\tau/\beta,
        \qquad
        v=\tau'/\beta,
        \qquad
        \theta_-=1-\frac{u+v}{2}.
    \]

    \item For each eigenmode \(a\), compute
    \[
        y_a=-x_a,
        \qquad
        c_a^-=\max(0,y_a).
    \]

    \item Use
    \[
        S_m^{-1}=S_m^\dagger.
    \]
    Build
    \[
        (M_m^-)_{ab}
        =
        \e^{-\theta_- c_a^-}
        (S_m^{-1})_{ab}
        \e^{-(1-\theta_-)c_b^-}
        +
        \delta_{ab}\e^{y_a-c_a^-}.
    \]

    \item Build
    \[
        (L_-)_{aa}
        =
        \exp\sq{uy_a-(1-\theta_-)c_a^-},
    \]
    \[
        (R_-)_{aa}
        =
        \exp\sq{(1-v)y_a-\theta_-c_a^-}.
    \]

    \item Solve
    \[
        M_m^-Y=R_-.
    \]

    \item Form
    \[
        K_m^< = -L_-Y.
    \]

    \item Transform back:
    \[
        \bar G_m^<(\tau,\tau')=VK_m^<V^T.
    \]
\end{enumerate}

\section{What This Stabilizes and What It Does Not}

The diagonal scaling removes artificial large numbers generated by
\(\e^{\pm\beta\Lambda}\).  In both branches:
\begin{itemize}[leftmargin=2em]
    \item the left prefactors are exponentials of non-positive numbers;
    \item the right prefactors are exponentials of non-positive numbers;
    \item the diagonal term in the middle matrix is at most \(1\);
    \item the mixing term is a unitary matrix element multiplied by two
    factors at most \(1\).
\end{itemize}
Thus no entry in the scaled construction contains a large exponential.

The method does not remove true singularities.  If \(D+S_m\) or
\(H+S_m^{-1}\) is close to singular, then the Green function itself is
ill-conditioned in that sector.  The scaled matrices \(M_m^\pm\) will then
also be ill-conditioned.  This is the correct numerical signal of a physical
or algebraic near-zero mode, not a failure of the scaling.

\section{Relation to the Current Direct Formula}

The direct implementation usually performs the following steps:
\[
    U(\beta)=\e^{-\beta\bh},
    \qquad
    B_m=\bI+\Omega_m^{-1}U(\beta),
    \qquad
    A_m=B_m^{-1},
\]
and later evaluates
\[
    \e^{-\tau\bh}A_m\e^{\tau'\bh}
    \quad\text{or}\quad
    \e^{-\tau\bh}(A_m-\bI)\e^{\tau'\bh}.
\]
This is correct, but it allows the large and small exponentials to appear
before cancellation.

The stable kernel method instead works directly with
\[
    K_m^> =
    L_+(M_m^+)^{-1}R_+,
    \qquad
    K_m^< =
    -L_-(M_m^-)^{-1}R_-,
\]
where all exponentials in \(L_\pm\), \(R_\pm\), and \(M_m^\pm\) have
non-positive exponents.  It is exactly equivalent to the original Green
function formula, up to ordinary roundoff and the conditioning of the scaled
linear solve.

\end{document}
